% Proved after full root and companion audits; exact candidate and outputs retained.
\subsection{A Frobenius-orbit collision at the next tame boundary}

Fix \(p\ge5\), put \(q=p^2\), \(N=q-1\), and set
\[
 e=(p+2)(p-1)=q+p-2,\qquad s=e(p+1)=(p+2)N.
 \tag{B-RD70.1}
\]
The complete mixed-characteristic \((0,p)\) field \(K\) has
algebraically closed residue field, \(\mu_p\subset K\), and
absolute ramification \(e\). Thus \(p\nmid e\),
\(q<e<2q\), and \(e<p^3\).
We use the actual good tame extension \(L/K\) of degree \(N\),
and the tangent exponent \(c\), exactly as in B69.
For a two-slope group its high line, parameter and valuations are
\[
 H,\qquad 0<B<pe,\quad B\equiv c\pmod{p+1},\qquad
 \alpha=s-B,\quad\beta=B/p .
 \tag{B-RD70.2}
\]
The conditional comparison rules and the actual high/low quotient
identities of Lemmas~\ref{lem:B-RD69-rules}
and~\ref{lem:B-RD69-quotients} apply at this \(e\): neither
lemma requires \(e<p^2\).

\begin{lemma}[Recovery when the high line is common]
\label{lem:B-RD70-high}
\status{Proved; complete independent root and companion proof and exact-code audits passed.}
Let two whole potentially supersingular height-two groups have
isomorphic actual \(p^3\)-torsion, two-slope first-level polygons,
and corresponding high lines. Their finite special inertia
characters agree.
\end{lemma}
\begin{proof}
The actual comparison identifies their finite determinant
characters. As in (B-RD69.10), their tangent exponents therefore
satisfy \(c_1\equiv c_2\pmod{p-1}\).
Their common high-line character gives
\[
 c_1-B_1\equiv c_2-B_2\pmod N .
 \tag{B-RD70.3}
\]
If either \(B_i\) is prime to \(p\), the common residual action
over \(L\) is wild, since its low coordinates have nonintegral
\(v_L\)-valuation \(B_i/p\). The first rule of
Lemma~\ref{lem:B-RD69-rules} then proves equality of the
characters or excludes a mixed-divisibility comparison.
This includes the possible critical value \(B_i=e\).

For the remaining cases both \(B_i\) are divisible by \(p\);
in particular neither equals \(e\).
If both \(B_i<e\), write \(B_i=pb_i\).
The common norm and (B-RD70.3) imply
\(b_1\equiv b_2\pmod{p-1}\), while
\[
 0<b_i<e/p=p+1-2/p,\qquad b_i\in\{1,\ldots,p\}.
 \tag{B-RD70.4}
\]
Consequently either \(b_1=b_2\) or
\(\{b_1,b_2\}=\{1,p\}\).
Equality gives \(c_1=c_2\pmod N\).
In the unequal case interchange the two models if necessary
so that \(b_1=1\), \(b_2=p\). Then
\[
 c_2-c_1\equiv p(p-1)\pmod N,\qquad
 c_1\equiv B_1=p\pmod{p+1}.
 \tag{B-RD70.5}
\]
Multiplying the latter congruence by \(p-1\) shows
\((p-1)c_1\equiv p(p-1)\pmod N\), so \(c_2\equiv pc_1\pmod N\).
The two unordered special orbits \(\{c_i,pc_i\}\) are equal.
Here unequal numerical parameters give the same character;
one must not incorrectly infer equality of the parameters.

If both \(B_i>e\), quotient by the common high line.
The inherited actual \(p^2\)-comparison has parameters
\(s-B_i\), both prime to \(p\), because \(p\nmid s\).
Lemma~\ref{lem:B-RD69-rules} recovers the quotient characters,
and isogeny invariance recovers the original ones.
If one \(B_i\) is below \(e\) and the other above, the high
quotient parameters are respectively \(pB_i\) and \(s-B_j\).
Their opposite \(p\)-divisibility contradicts the same lemma.
These are all possibilities.
\end{proof}

\begin{lemma}[A constraint on exchanged tame high lines]
\label{lem:B-RD70-exchange}
\status{Proved; complete independent root and companion proof and exact-code audits passed.}
Suppose two two-slope groups have isomorphic residual modules
with tame inertia. If the isomorphism does not identify their
high lines, then
\[
                         p-1\mid B_1+B_2.
 \tag{B-RD70.6}
\]
In particular two such groups with \(B_1=B_2=p^2\) must have
corresponding high lines.
\end{lemma}
\begin{proof}
The high eigenvalue of the tame generator is
\(\zeta^{c_i-B_i}\). Its determinant is
\(\zeta^{(p+1)c_i}\), because \(\mu_p\subset K\).
Thus the other residual eigenvalue is \(\zeta^{pc_i+B_i}\).
These are characters of the actual residual representation,
not of the special isocrystal.
Tame residual inertia is semisimple. If its two eigenvalues
are distinct, every stable line other than its high line
is its other eigenspace. If the representation is scalar,
the other eigenvalue has the same value and the following
identities remain valid for every line. Hence distinct
corresponding high lines force
\[
 c_1-B_1\equiv pc_2+B_2,\qquad
 c_2-B_2\equiv pc_1+B_1\pmod N .
 \tag{B-RD70.7}
\]
Multiplying the first congruence by \(p\), adding the second,
and using \(p^2\equiv1\pmod N\), gives
\(N\mid(p+1)(B_1+B_2)\). This is precisely (B-RD70.6).
For \(B_1=B_2=p^2\), it would require \(p-1\mid2p^2\),
or \(p-1\mid2\), impossible for \(p\ge5\).
\end{proof}

\begin{proposition}[Tame reducible comparisons at the next boundary]
\label{prop:B-RD70-tame}
\status{Proved; complete independent root and companion proof and exact-code audits passed.}
Two whole potentially supersingular height-two groups at the
index (B-RD70.1), with two-slope first-level polygons, tame
residual inertia and isomorphic actual \(p^3\)-torsion, have
equal finite special inertia characters.
\end{proposition}
\begin{proof}
Their first-level modules are rational over \(L\), so write
\(B_i=pb_i\), with \(0<b_i<e\).
If both \(b_i\) are prime to \(p\), the second comparison
rule of Lemma~\ref{lem:B-RD69-rules} already recovers their
characters from actual \(p^2\)-torsion.
Suppose exactly one is prime to \(p\), say \(b_1=b\), while
\(b_2=pk\). Choose corresponding residual vectors outside
both high lines and corresponding primitive lifts.
The first lift has valuation \(b/p^2\). Its denominator
forces the common point field to have degree \(p^2\)
over \(L\); all translations by the \(L\)-rational first-level
module occur. Its normalized coordinate valuations in that
field are \(b\) and \(pk\).

If the high lines differ, choose a translation high in the
first group and low in the second. The first coordinate has
prime-to-\(p\) valuation, so the actual lower break is
\(p^2\alpha_1-b\). The lower bound for the second gives
\[
 p^2(\alpha_1-pk)\le b-pk .
 \tag{B-RD70.8}
\]
But \(\alpha_1=s-pb\ge e+p\), and \(pk<e\).
Consequently the left side is greater than \(p^3\),
whereas the right side is less than \(e<p^3\).
Thus the high lines agree, and
Lemma~\ref{lem:B-RD70-high} proves the result.
Unlike the earlier proof below \(p^2\), this step does not
assert \(p^2k>e\); it allows the new case \(k=1\).

It remains that \(b_i=pk_i\) for both groups, so
\(B_i=p^2k_i\). If the high lines agree, use the same lemma.
Suppose they do not. If either \(B_i>e\), label that group
first, quotient it by its high line, and quotient the other
by the corresponding low line. The lines are inertia stable
because they correspond under the actual Galois-module
isomorphism. The inherited actual \(p^2\)-comparison has
parameters
\[
                  s-B_1\quad\text{and}\quad B_2/p=pk_2.
 \tag{B-RD70.9}
\]
The first is prime to \(p\), the second divisible by \(p\);
Lemma~\ref{lem:B-RD69-rules} forbids this comparison.
Neither original \(B_i\) equals \(e\), since \(p\mid B_i\)
and \(p\nmid e\). The only remaining possibility is
\(B_1<e\) and \(B_2<e\). Since \(p^2<e<2p^2\), this
forces \(B_1=B_2=p^2\).
Lemma~\ref{lem:B-RD70-exchange} now says that their high
lines must agree, contradicting the supposition.
This exhausts every tame divisibility case.
\end{proof}

\begin{proposition}[All whole formal comparisons at this index]
\label{prop:B-RD70-formal}
\status{Proved; complete independent root and companion proof and exact-code audits passed.}
At \(e=(p+2)(p-1)\), actual \(p^3\)-torsion determines the
finite special inertia character among whole potentially
supersingular height-two groups.
\end{proposition}
\begin{proof}
Proposition~\ref{prop:B-RD69-one-slope}, which has no upper
bound on \(e\), excludes a mixture of one-slope and two-slope
polygons and recovers two one-slope characters.
For two two-slope groups with nontrivial residual wild inertia,
the high line is the unique wild fixed line, hence common;
Lemma~\ref{lem:B-RD70-high} applies.
Tame residual inertia is precisely the preceding proposition.
The argument covers all finite determinant characters without
a square-class normalization.
\end{proof}

\begin{theorem}[The next tame family has exact index two]
\label{thm:B-RD70}
\status{Proved; complete independent root and companion proof and exact-code audits passed.}
For every prime \(p\ge5\), in the full B.4 monodromy setting
with height two, perfect residue field and \(\mu_p\subset K\),
\[
 \boxed{\displaystyle
       \nu_{\min}\bigl(p,(p+2)(p-1),2\bigr)=2 .}
 \tag{B-RD70.10}
\]
In particular \(\nu_{\min}(5,28,2)=2\).
\end{theorem}
\begin{proof}
For the upper bound pass to the completed maximal unramified
base. Since \(p\nmid e\), the cyclotomic image modulo \(p^3\)
has order \(p^2\); the determinant argument of
Theorem~\ref{thm:B-RD69} identifies the total dimension
\(d\in\{0,1,2\}\) and the finite determinant.
Dimensions zero and two, and two systems filtered by
potentially good height-one pieces, are detected residually
as in that proof. In dimension one a mixed whole/filtered
comparison is excluded by
Proposition~\ref{prop:B-RD63-mixed}, since \(e<p^3\).
Two whole potentially supersingular groups are covered by
Proposition~\ref{prop:B-RD70-formal}.
Thus actual \(p^3\)-torsion always suffices.
As \(e>p\), Theorem~\ref{thm:B-RD63-arc} with \(n=2\)
provides permitted height-two systems with equal actual
\(p^2\)-torsion and different requested characters.
Index one therefore fails and index two is least.
\end{proof}

\begin{remark}[The orbit collision and the next limitation]
\label{rem:B-RD70-scope}
\status{Proved; complete independent root and companion proof and exact-code audits passed.}
The new small-parameter collision is \(\{b_1,b_2\}=\{1,p\}\).
It preserves the unordered special character by exchanging its
two Frobenius-conjugate eigenvalues. The proof does not rule out
all unequal Newton parameters.
The new tame case \(B_1=B_2=p^2\) is handled by actual residual
eigenvalues, not by a square-determinant quotient-thickness bound.
At \(e=(p+3)(p-1)\), the small interval first admits
\(\{2,p+1\}\), which is not automatically a Frobenius orbit;
no conclusion for that index follows from this proof.
No absolute ramification divisible by \(p\) is treated.
The retained exact supplement
\path{code/reduction_descent/verify_next_tame_boundary.py}
checks twelve primes from \(5\) through \(43\).
It verifies \(8{,}496\) common-high-line character orbits,
including \(528\) unequal-parameter Frobenius collisions,
\(7{,}704\) critical exchanged-line pairs,
\(280{,}368\) mixed point-field inequalities, and every
high/low quotient-divisibility case in the stated finite range.
The first unexcluded later arithmetic pattern
\((b_1,b_2)=(2,p+1)\) is retained explicitly; it is not
asserted to arise from formal groups with isomorphic torsion.
The author run passed in \(0.14006940000399482\) seconds at
21:40:22 UTC. An unused malformed draft block was removed
before the first execution, and that unexecuted draft is
preserved. No executed numerical check failed.
These finite checks supplement the preceding all-prime proofs.
The complete independent root and companion audits passed,
with author-output-preserving exact reruns in
\(0.16807877800602\) and \(0.1444874750013696\) seconds,
respectively. Both full audit records and the exact
prepromotion candidate are retained.
\end{remark}
